\clearpage
\section{선별 문제 풀이 II}
\label{sec:norm-trace-solutions-b}

\subsection*{문제 \thechapter.17: 노름과 대각합의 탑 공식}

\begin{solution}
대수적 폐포 $\overline K$를 고정하자. $E/K$의 서로 다른 매장을
\[
  \sigma_1,\dots,\sigma_r:E\hookrightarrow\overline K
\]
라 하고, 각 $\sigma_i$를 $\overline K$의 자동동형
$\widetilde\sigma_i$로 연장한다. $L/E$의 서로 다른 매장을
\[
  \tau_1,\dots,\tau_s:L\hookrightarrow\overline K
\]
라 하면
\[
  \widetilde\sigma_i\circ\tau_j
  \qquad(1\le i\le r,\ 1\le j\le s)
\]
가 $L/K$의 서로 다른 모든 매장을 준다.

비분리차수를
\[
  q_1=[E:K]_{\mathrm i},
  \qquad
  q_2=[L:E]_{\mathrm i}
\]
라 하면
\[
  [L:K]_{\mathrm i}=q_1q_2.
\]
따라서 체매장 공식으로
\[
\begin{aligned}
  \Tr_{L/K}(a)
  &=q_1q_2\sum_{i,j}
    \widetilde\sigma_i(\tau_j(a))\\
  &=q_1\sum_i
    \widetilde\sigma_i
    \left(q_2\sum_j\tau_j(a)\right)\\
  &=\Tr_{E/K}\bigl(\Tr_{L/E}(a)\bigr).
\end{aligned}
\]
또한
\[
  \Nm_{L/E}(a)
  =\left(\prod_j\tau_j(a)\right)^{q_2}
\]
이므로
\[
\begin{aligned}
  \Nm_{E/K}\bigl(\Nm_{L/E}(a)\bigr)
  &=\left(
    \prod_i
    \widetilde\sigma_i
    \left(
      \left(\prod_j\tau_j(a)\right)^{q_2}
    \right)
    \right)^{q_1}\\
  &=\left(
    \prod_{i,j}
    \widetilde\sigma_i(\tau_j(a))
    \right)^{q_1q_2}\\
  &=\Nm_{L/K}(a).
\end{aligned}
\]
따라서
\[
  \boxed{
  \Tr_{L/K}
  =\Tr_{E/K}\circ\Tr_{L/E}},
  \qquad
  \boxed{
  \Nm_{L/K}
  =\Nm_{E/K}\circ\Nm_{L/E}}.
\]
\end{solution}

\subsection*{문제 \thechapter.19: 분리성과 대각합 쌍}

\begin{solution}
(a)$\Rightarrow$(b): $L/K$가 분리이면
\[
  \Tr_{L/K}
  =\sigma_1+\cdots+\sigma_n
\]
이다. 서로 다른 체매장은 선형독립이므로 이 합은 영사상이 아니다.

(b)$\Rightarrow$(c): $x\in L$가 모든 $y\in L$에 대해
\[
  \Tr_{L/K}(xy)=0
\]
을 만족한다고 하자. $x\ne0$이면 $xL=L$이므로
$\Tr_{L/K}$가 $L$ 전체에서 영이 되어 (b)에 모순이다.
따라서 $x=0$이고 대각합 쌍은 비퇴화이다.

(c)$\Rightarrow$(d): 기저 $x_1,\dots,x_n$에서 대각합 쌍의
그람행렬이
\[
  G=(\Tr(x_ix_j))
\]
이다. 쌍이 비퇴화이면 $G$는 가역이고 $\det G\ne0$이다.

(d)$\Rightarrow$(a): 만약 $L/K$가 비분리이면 비분리차수는
표수 $p$의 비자명한 거듭제곱이고
\[
  \Tr_{L/K}=0.
\]
그러면 모든 그람행렬이 영행렬이 되어 판별식이 $0$이다.
(d)에 모순이므로 $L/K$는 분리이다.
\end{solution}

\subsection*{문제 \thechapter.21: 유한체 노름의 전사성}

\begin{solution}
$\F_{q^n}^\times$는 위수 $q^n-1$인 순환군이다.
생성원을 $g$라 하자. 노름 공식은
\[
  \Nm(g)
  =g^{1+q+\cdots+q^{n-1}}
  =g^{(q^n-1)/(q-1)}.
\]
순환군에서 $g^e$의 위수는
\[
  \frac{q^n-1}{\gcd(q^n-1,e)}.
\]
여기서 $e=(q^n-1)/(q-1)$이므로
\[
  \operatorname{ord}(\Nm(g))=q-1.
\]
따라서 $\Nm(g)$는 $\F_q^\times$의 생성원이고 노름은 전사이다.

제1동형정리에서
\[
  \left|\ker\Nm\right|
  =\frac{|\F_{q^n}^\times|}{|\F_q^\times|}
  =\boxed{\frac{q^n-1}{q-1}}.
\]
\end{solution}

\subsection*{문제 \thechapter.24: 서로소 차수의 근 하강}

\begin{solution}
$b\in L^\times$가 $b^n=a$를 만족한다고 하자. 노름의 곱셈성과
$a\in K$에 대한 스칼라 공식에서
\[
  \Nm_{L/K}(b)^n
  =\Nm_{L/K}(b^n)
  =\Nm_{L/K}(a)
  =a^m.
\]
$\gcd(m,n)=1$이므로 어떤 정수 $r,s$가 존재하여
\[
  rm+sn=1.
\]
이제
\[
  c=\Nm_{L/K}(b)^r a^s\in K^\times
\]
라 두면
\[
\begin{aligned}
  c^n
  &=\Nm_{L/K}(b)^{rn}a^{sn}\\
  &=a^{rm}a^{sn}\\
  &=a^{rm+sn}=a.
\end{aligned}
\]
따라서 $a=c^n$이고 $a$는 이미 $K$에서 $n$제곱이다.
음의 지수는 $a,b\ne0$이므로 문제가 없다.
\end{solution}

\subsection*{문제 \thechapter.26: 종결식과 노름}

\begin{solution}
$A=K[X]/(f)$에서 $\alpha=X+(f)$라 하자. $f$가 기약이므로
$A=L$은 체이고, 기저
\[
  1,\alpha,\dots,\alpha^{n-1}
\]
를 갖는다. Chapter 2의 종결식 해석에 따르면
\[
  \Res(f,g)
\]
는 $A$에서 $g(\alpha)$를 곱하는 $K$-선형사상의 행렬식이다.
그러나 이는 정의상
\[
  \Nm_{L/K}(g(\alpha))
\]
이다. 따라서
\[
  \boxed{
  \Nm_{L/K}(g(\alpha))=\Res(f,g)}.
\]

특히 $g(X)=x-X$를 택하면
\[
  g(\alpha)=x-\alpha.
\]
$f$가 일계수이므로
\[
  \Res(f,x-X)=f(x),
\]
또는 특성다항식의 정의에서 직접
\[
  \Nm_{L/K}(x-\alpha)
  =\det(xI-m_\alpha)
  =\chi_{\alpha,L/K}(x)
  =f(x)
\]
를 얻는다.
\end{solution}

\subsection*{문제 \thechapter.29: 두 기저의 판별식}

\begin{solution}
첫 기저 $\mathcal B_1=(1,\sqrt5)$에 대해
\[
  \Tr(1)=2,\quad
  \Tr(\sqrt5)=0,\quad
  \Tr(5)=10.
\]
따라서
\[
  G_1=
  \begin{pmatrix}
    2&0\\
    0&10
  \end{pmatrix},
  \qquad
  D(\mathcal B_1)=20.
\]

$\omega=(1+\sqrt5)/2$라 하자. $\omega$의 최소다항식은
\[
  X^2-X-1
\]
이므로
\[
  \Tr(\omega)=1,\qquad
  \Nm(\omega)=-1.
\]
또한 $\omega^2=\omega+1$이므로
\[
  \Tr(\omega^2)=\Tr(\omega)+\Tr(1)=1+2=3.
\]
따라서
\[
  G_2=
  \begin{pmatrix}
    2&1\\
    1&3
  \end{pmatrix},
  \qquad
  D(\mathcal B_2)=6-1=5.
\]

두 기저는
\[
  (1,\omega)
  =(1,\sqrt5)
  \begin{pmatrix}
    1&1/2\\
    0&1/2
  \end{pmatrix}
\]
로 연결되고 변환행렬의 행렬식은 $1/2$이다. 따라서
\[
  D(\mathcal B_2)
  =\left(\frac12\right)^2D(\mathcal B_1)
  =\frac14\cdot20=5
\]
로 기저변환 공식과 일치한다.
\end{solution}

\subsection*{문제 \thechapter.30: $1-\zeta_p$의 노름과 대각합}

\begin{solution}
$\zeta_p$의 최소다항식은
\[
  \Phi_p(X)=1+X+\cdots+X^{p-1}
\]
이다. 최소다항식의 노름 표현에서
\[
  \Nm_{\Q(\zeta_p)/\Q}(1-\zeta_p)
  =\Phi_p(1)=p.
\]
따라서
\[
  \boxed{
  \Nm_{\Q(\zeta_p)/\Q}(1-\zeta_p)=p.}
\]

또한
\[
  \Tr_{\Q(\zeta_p)/\Q}(\zeta_p)=-1
\]
이고 확장차수는 $p-1$이므로
\[
\begin{aligned}
  \Tr_{\Q(\zeta_p)/\Q}(1-\zeta_p)
  &=\Tr(1)-\Tr(\zeta_p)\\
  &=(p-1)-(-1)=p.
\end{aligned}
\]
따라서
\[
  \boxed{
  \Tr_{\Q(\zeta_p)/\Q}(1-\zeta_p)=p.}
\]
\end{solution}
